\section{Case Study: Closing the Slow Corner}
\label{sec:closure}

\begin{table*}[t]
\centering\small
\caption{Engine closure iterations. Setup slack (worst endpoint) is in ns at a \SI{150}{ns} clock with extracted parasitics. Energy uses activity-annotated routed power at TT. Completed routed runs pass routing, Magic DRC, main KLayout DRC and LVS checks; fix 1 stopped before completion. Electrical-limit violations remain. Antenna counts are open nets after repair.}
\label{tab:closure}
\setlength{\tabcolsep}{4.5pt}
\begin{tabular}{@{}lrrrrrrl@{}}
\toprule
Run & GRT ovf. & TT & SS nom & SS max & Antenna & pJ/MAC & What changed\\
\midrule
First route (var.\ C) & 575 & $-72.9$ & $-251$ & $-307$ & 130 & 241 & local partials; no post-placement repair\\
Closure ECO 1 & -- & $+55.0$ & $-40.0$ & $-55.5$ & 4 & 256 & post-route repair; antenna diodes\\
Closure ECO 2 & -- & $+56.7$ & $-36.2$ & $-53.1$ & 4 & -- & more aggressive repair\\
Fix 1 (stopped) & 3,206 & -- & -- & -- & -- & -- & register before FP32 stage\\
Fix 2 & 9 & $+67.6$ & $+21.6$ & $+8.3$ & 11 & 287 & $+$ split 320:1 select\\
\textbf{Fix 2, final} & \textbf{9} & $\mathbf{+67.5}$ & $\mathbf{+21.4}$ & $\mathbf{+8.1}$ & \textbf{0} & \textbf{287} & $+$ 10 antenna-repair rounds\\
\bottomrule
\end{tabular}
\end{table*}

Timing closure of the $1280\times320$ engine illustrates both how the open flow behaves at scale and why architectural fixes beat tool tuning. \Cref{tab:closure} summarises each iteration. Rows with timing entries use extracted-parasitic (RCX/SPEF) static timing at a \SI{150}{ns} clock (\SI{6.67}{MHz}).

\paragraph{Congestion first.}
The first full placements could not be routed. Global routing reported 3.37\,M overflowed resources, because one set of radix-16 partial products fed all 320 lanes across the die. Giving each 32-lane group its own cycle-identical copy of the partial-product registers (``local partials'') cut overflow to 260\,k, and then to 575 once post-placement repair was disabled. Detailed routing then converged to zero violations.

\paragraph{Tool-level closure plateaus.}
With routing clean, the typical corner (TT, \SI{25}{\celsius}, \SI{3.3}{V}) met timing with $+55$\,ns of slack. The slow corner (SS, \SI{125}{\celsius}, \SI{3.0}{V}) failed by $-40$\,ns (nominal RC) and $-55.5$\,ns (worst-case RC). We ran two engineering-change (ECO) passes, re-entering the flow after global routing. They used more aggressive repair: a \SI{45}{ns} setup-slack margin, buffering allowed up to 50\% of cells, 100\% of failing endpoints targeted, and 20 antenna-repair iterations. These passes recovered only about 4\,ns (to $-36.2$/$-53.1$\,ns), and the \emph{same} 30 endpoints failed in both. We stopped tuning there: the failing paths were too deep for one cycle at SS, whatever the repair settings.

\paragraph{Diagnosis.}
All 30 failing paths start at the column register. They then pass through:
\begin{itemize}
\item the 320:1 accumulator select;
\item the zero-point lookup and its correction multiply and subtract;
\item the integer-to-floating-point conversion.
\end{itemize}
They end in the first register stage of the FP32 affine block.

The two first-order models disagreed on the size of the problem. Using estimated parasitics after global routing, the pre-fix design showed $+36$\,ns of SS slack. Its routed result was $-36$\,ns, a 72\,ns gap. The gap comes from the die-wide 320:1 select: its fan-in wires are long and only appear once detailed routing is done.

\paragraph{Fix~1: one elastic register.}
We first inserted an elastic valid/ready register between the integer engine output and the FP32 affine stage; weight scale, bias and column index move through it with the data. Throughput is unchanged, and latency grows by one cycle (4,163 to 4,164 cycles per projection). The 54-vector bench still matched all 17,280 FP32 outputs bit-exactly. Estimated SS slack on the column path rose from $+36$ to $+61$\,ns. Applying the 72\,ns gap leaves an expected routed slack of about $-10$\,ns, still failing. Global-route overflow also rose to 3,206. We therefore stopped this run at detailed routing.

\paragraph{Fix~2: split the select with look-ahead group registers.}
Each 32-lane group now registers its own lane for the \emph{look-ahead} column, so the output path sees only a 10:1 select of group registers. The look-ahead column is $c{+}1$ on an output handshake and $c$ otherwise. The zero-point correction is registered the same way. A single \texttt{PREP} state between the last accumulation and the first output fills these registers. The change is generated by a script from the original RTL. It adds one more cycle (4,165 per projection), keeps one output per cycle, and again passes all 54 vectors bit-exactly.

The physical effect was large. Global-route overflow fell to 9, and detailed routing converged in about 75 minutes. The completed geometric/equivalence checks report zero routing, Magic and main KLayout DRC violations, a unique Netgen LVS match, and zero GDS XOR differences; 11 antenna nets remain open, against 4 before. All nine corners (TT/SS/FF $\times$ min/nom/max RC) meet setup and hold at \SI{150}{ns}. The worst case is $+8.34$\,ns at SS with worst-case RC, against $-53.1$\,ns before.

\paragraph{Antenna close-out.}
The 11 open antenna nets were not a design problem. The router's post-route antenna repair had stopped at its default three rounds while still converging (199, 41, 24, then 11 violating nets). Re-running from the same detailed-routing input with up to ten rounds, and nothing else changed, took the count to zero: 8, 5, 5, 5, 1, 1, then 0, adding about 1,000 diodes (7,452 in total). The listed open-flow checks were then repeated: routing, Magic and KLayout DRC, XOR, disconnected-pin and power-grid checks are all zero, and LVS again matches uniquely. Timing moved by less than \SI{0.3}{ns} in every corner (worst $+8.10$\,ns setup, $+0.32$\,ns hold), and routed power with the same activity trace is unchanged. The run took 6.7 hours on a laptop, two of them in KLayout DRC.


\paragraph{Cost of closure.}
Routed, activity-annotated power for fix~2 is \SI{0.188}{W} at \SI{150}{ns} at TT. That is \SI{117}{\micro\joule} per projection, or \SI{287}{pJ} per MAC, about 12\% more than the ECO-1 design (\SI{256}{pJ}/MAC). The increase accompanies extra group registers, registered correction and repair buffering; the flattened report cannot assign a measured energy increment to each change. Both power runs annotate about 68\% of pins from a zero-delay gate-level trace of one real-speech projection. The unannotated pins belong to repair, antenna and diode cells that do not exist in the synthesized netlist.

\paragraph{Lesson.}
Two passes of repair-setting tuning, about 17 hours of compute, did not move the failing endpoints. One RTL change of under 40 lines, verified bit-exactly on the full bench within minutes, closed every corner and also removed the routing congestion. The general rule we draw is to pipeline wide selects and long fan-in paths at RTL before physical design. A small set of path-level timing estimates, calibrated against one fully routed run, is enough to decide when to stop tuning.
